\section{Strumenti}
\label{suc:Strumenti}

Nella seguente sezione sono elencati gli strumenti utilizzati dal team per lo svolgimento delle attività di progetto.

\subsection{Discord}
\label{subsection:discord}
Link - \href{https://discord.com/}{https://discord.com/} \\
Il team utilizza la piattaforma Discord per le riunioni interne e per condividere materiale utile allo svolgimento del progetto

\subsection{Telegram}
\label{subsection:telegram}
Link - \href{https://web.telegram.org/}{https://web.telegram.org/} \\
Il team utilizza la piattaforma Telegram per la messaggistica istantanea e viene utilizzata principalmente 
per condividere informazioni o in caso di dubbi 

\subsection{Latex}
\label{subsection:latex}
Link - \href{https://www.latex-project.org/}{https://www.latex-project.org/} \\
Il gruppo utilizza Latex per redigere tutti i documenti ufficiali presenti nel \glossario{repository} 

\subsection{Git}
\label{subsection:git}
Link - \href{https://git-scm.com/}{https://git-scm.com/} \\
Il gruppo ha scelto di utilizzare Git come version control system

\subsection{GitHub}
\label{subsection:github}
Link - \href{https://github.com/}{https://github.com/} \\
Il gruppo ha deciso di utilizzare GitHub come piattaforma di hosting per lo sviluppo software collaborativo. GitHub
offre controllo di versione, \glossario{repository} Git, gestione del progetto e automazioni.

\subsection{Google Docs}
\label{subsection:google_docs}
Link - \href{https://docs.google.com/}{https://docs.google.com/} \\
Il gruppo ha deciso di utilizzare Google Docs per scrivere documenti non ufficiali online in modo collaborativo in tempo reale.

\subsection{Gmail}
\label{subsection:gmail}
Link - \href{https://mail.google.com/}{https://mail.google.com/} \\
Il gruppo ha scelto di utilizzare Gmail come servizio di posta elettronica per la comunicazione ufficiale con i 
professori e l'azienda proponente.

\subsection{Google Sheets}
\label{subsection:google_sheets}
Link - \href{https://docs.google.com/spreadsheets/}{https://docs.google.com/spreadsheets/} \\
Il gruppo ha deciso di utilizzare Google Sheets sia per scrivere in fogli di calcolo non ufficiali online in modo 
collaborativo in tempo reale sia per organizzare dati in modo strutturato e automatico attraverso le github actions.

\subsection{Google Slides}
\label{subsection:google_slides}
Link - \href{https://docs.google.com/presentation/}{https://docs.google.com/presentation/} \\
Il gruppo ha deciso di utilizzare Google Slides per la creazione dei diari di bordo e per le presentazioni relative 
alle revisioni di avanzamento.

\subsection{StarUML}
\label{subsection:staruml}
Link - \href{https://staruml.io/}{https://staruml.io/} \\
Il gruppo ha deciso di utilizzare StarUML per la creazione di tutti i diagrammi UML presenti nella documentazione.

\subsection{Visual Studio Code}
\label{subsection:visual_studio_code}
Link - \href{https://code.visualstudio.com/}{https://code.visualstudio.com/} \\
Il gruppo ha deciso di utilizzare Visual Studio Code come editor sia per la documentazione che per la codifica.

\subsection{LTeX+}
\label{subsection:ltex+}
Link - \href{https://ltex-plus.github.io/ltex-plus/}{https://ltex-plus.github.io/ltex-plus/} \\
Il gruppo ha deciso di utilizzare LTeX+ per il controllo grammaticale dei documenti 

